% !TeX root = ../drafts/06-bus-interactions.tex
\section{Bus interactions}\label{sec:omc}

Domain-separated bus channels bind CPU transitions, immutable ELF bytes, RAM events, register events, sorted histories, and witness-copy iterations. Every coordinate is a degree-at-most-two polynomial over $\K$. Packed addresses and values retain their integer bit representations; field multiplication of a lookup counter does not change the machine's addressing convention.

\subsection{CPU state}\label{sec:state}

An enabled CPU row pulls $\tup{\dsep{CPU},\pc,\rvcycle,0}$ and pushes its successor. Ordinary successors are $\tup{\dsep{CPU},\nxt{\pc},\rvcycle+1,0}$; EXIT pushes $\tup{\dsep{CPU},0,\rvcycle+1,1}$. One public boundary pushes $\tup{\dsep{CPU},\rventry,0,0}$ and pulls $\tup{\dsep{CPU},0,\rvcycles,1}$. Enabled cycles lie in $[0,\rvcycles)$, with $0<\rvcycles<2^{32}$, and increment without wrap.

\begin{proposition}[Balanced transition rows]\label{prop:walk}
Let a finite multiset of rows pull states $\bpull(r)$ and push states $\bpush(r)$. If adding one initial push and one final pull balances the state channel, its rows decompose into a walk from the initial state to the final state and closed walks.
\end{proposition}
\begin{proof}
Make one directed edge per row, from its pulled state to its pushed state, and add an edge from the final state to the initial state. Balance makes every vertex's in-degree equal its out-degree. Decompose this finite balanced multigraph into edge-disjoint closed walks. Removing the added edge leaves the asserted initial-to-final walk.
\end{proof}
For enabled RV64IM rows, a closed walk is impossible: the bounded integer cycle strictly increases. Thus all enabled rows belong to the execution. Disabled padding has matching null bus tuples, rather than extra architectural cycles (\S\ref{sec:e2e-pad}).

\subsection{Immutable lookup protocol}\label{sec:lookup}

The counted lookup argument is used only for the public ROM, not for mutable RAM. Its algebraic core follows offline memory checking~\cite{BEGKN94} and \cite[Protocol~4.27]{DP23}. A public array has entries $\tup{a_i,\lut[i]}$, where $a_i$ is an ordinary integer address packed into $\K$; its value may have several coordinates. The following generic argument does not require exponent-encoded addresses.
\begin{itemize}[itemsep=1pt,topsep=3pt]
\item \textbf{Seed:} for every array entry, push $\tup{\sep,a_i,1,\lut[i]}$.
\item \textbf{Read:} pull $\tup{\sep,a,c,v}$ and push $\tup{\sep,a,\gen c,v}$.
\item \textbf{Finalize:} for every entry, pull $\tup{\sep,a_i,\cntfin[i],\lut[i]}$.
\end{itemize}
If an entry is read $A_i$ times, the honest finalize count is $\gen^{A_i}$ and its successive read counts start at $1$. Only the counts, not the addresses, use powers of $\gen$.

\begin{lemma}[Invariant count multisets]\label{lem:invcnt}
Let $\rcnts$ be a finite multiset over $\K$ with $\gen\rcnts=\rcnts$. If $|\rcnts|<2^{64}-1$, every element of $\rcnts$ is zero.
\end{lemma}
\begin{proof}
Multiplication by $\gen$ permutes the multiset and acts transitively on $\K^\times$. Every nonzero field element therefore has the same multiplicity. A nonzero multiplicity would contribute at least $2^{64}-1$ entries, contrary to the bound.
\end{proof}

\begin{theorem}[Lookup correctness]\label{thm:omc}
Suppose these are the only flushes with separator $\sep$, the bus balances, every read count is nonzero, and there are fewer than $2^{64}-1$ reads. Then every read pair $(a,v)$ equals some public entry $(a_i,\lut[i])$.
\end{theorem}
\begin{proof}
Fix a pair absent from the public array. Filtering the balanced bus by separator and this pair removes every seed and finalization. Its read-count multiset satisfies $\gen\rcnts=\rcnts$, so Lemma~\ref{lem:invcnt} makes every surviving count zero. Nonzero read counts force the multiset to be empty.
\end{proof}
\begin{corollary}[Addresses are in range]\label{cor:addrrange}
Under the same hypotheses, an enabled ROM read has a tagged key occurring in the public ROM. Instruction entries exist only for executable words, and the CPU checks program-counter alignment and bounds. RAM bounds are checked by the memory circuit.
\end{corollary}

\paragraph{Disabled reads and count product.} Let $e$ be a Boolean lookup-enable bit. Disabled reads emit the same tuple on both sides: separator, address, and value are multiplied by $e$, while the pull count remains $c$ and the push count is $c+e(\gen+1)c$. Thus inactive counts remain nonzero, rather than becoming zero. The third GKR product proves that all lookup-count columns have nonzero product. Enabled reads obey the theorem after the identical inactive tuples cancel. The fixed lookup rules have one instruction-word lookup per CPU row and one initialization lookup per RAM row; table-height caps keep the total below $2^{64}-1$. Finalize counts need no separate semantic check: an incorrect one can only unbalance its public entry.

\subsection{Mutable RAM and register histories}\label{sec:memchan}

Each architectural access emits an event $\tup{\dsep{RAM},\addr,\rvtime,v_{\mathrm{before}},v_{\mathrm{after}},w}$, or the corresponding \texttt{REG} tuple. The write indicator $w$ is Boolean. RAM events have $32$-bit addresses and $8$-bit values; register events have $5$-bit addresses and $64$-bit values. A read preserves its value. CPU and ECALL tables emit the events; a separate history table consumes exactly the same multiset.

RAM timestamps concatenate a $32$-bit cycle and $32$-bit slot, $\rvtime=\rvcycle\,2^{32}+\mathit{slot}$. Register timestamps concatenate the cycle and a $4$-bit slot, $\rvtime=\rvcycle\,2^4+\mathit{slot}$. These are bounded integer concatenations, not field multiplication. Register reads precede destination writes. Within a BLAKE2s ECALL, all input-byte events precede output-byte events. F192 multiplication has only register events. The witness-copy byte index supplies its RAM slot.

\paragraph{Sorted chain.} The history table has a row index, previous address/time/value, and the current event. One boundary seeds index zero with zero previous state; one terminal pull binds the padded row count and the announced final address/time/value. Every row advances the index by one without wrap. An event is active exactly when its index is below the announced real event count. Active rows are sorted lexicographically by address then time, with strictly increasing times for the same address. Equal-address successors require the current before-value to equal the preceding after-value. The unique seed and bounded, incrementing index exclude disconnected cycles, duplicate indices, and invented initial histories. Padding preserves the last history state while advancing only the row index.

\paragraph{Initialization and permissions.} At the first event for each address, RAM's before-value is the corresponding ELF file byte, proved by a ROM lookup, or zero for BSS and otherwise unmapped RAM. Segment read/write permissions constrain each active data event. Registers initialize to zero except \texttt{sp}$=2^{32}$; every access to \texttt{x0} has before- and after-value zero. Address, time, byte, and register ranges are enforced by bit decomposition. Therefore a prover cannot reset a mutable location to another advised value, substitute a host snapshot, or reorder overlapping writes without breaking a constraint or a bus equality.

\subsection{Public ELF ROM}\label{sec:bytecode}

The public ROM contains every file-backed byte at its virtual address shifted left once, and executable instruction words at aligned program counters shifted left once with the low bit set. Entries are sorted by these tagged keys. Every included instruction word has four executable bytes; a BSS byte contributes zero. Words wholly in BSS decode to no supported instruction and need not be materialized. A CPU row authenticates its raw $32$-bit instruction with one word lookup. First RAM accesses to file-backed data use the even-key byte entries. Writable data can subsequently change through RAM events; executable segments cannot be written.

\subsection{Public address columns}\label{sec:idxcol}

The ROM is padded to a power of two with sentinel key $2^{64}-1$ and zero value. Its key and value multilinears are evaluated from the public arrays,
\[
  \mle{\lut}(z)=\sum_{i\in\cube{\kbc}}\eq(z,i)\lut[i].
\]
Actual tagged fetch and RAM keys cannot equal the sentinel. The execution verifier evaluates the public program itself. In recursive composition this evaluation is constrained by the field-native circuit, whose program templates come from the trusted expected ELF, rather than executed by the RV64IM aggregation guest. No generator-power address polynomial or unresolved bytecode opening is part of the machine.
